\begin{answer}

    1) The projection of a vector $x^{(i)}$ onto a unit vector $u$ is by taking the dot product of $x^{(i)}$ and $u$ (which gives the length of the projection) and multiplying it by the direction vector $u$:

    $$f_u(x^{(i)}) = ((x^{(i)})^T u) u$$
    
    We can now use this to put into the objective function for which we are trying to find the most minimizing $u$:
    
    $$|x^{(i)} - f_u(x^{(i)})|^2_2 = |x^{(i)} - ((x^{(i)})^T u) u|^2_2$$
    
    With $|v|^2_2 = v^T v$, we can expand to:
    
    $$\left( x^{(i)} - ((x^{(i)})^T u) u \right)^T \left( x^{(i)} - ((x^{(i)})^T u) u \right)$$
    
    With the formula $(a-b)(a-b) = a^2 - 2ab + b^2$, we can expand this to:
    
    $$|x^{(i)} - f_u(x^{(i)})|^2_2 = \underbrace{|x^{(i)}|^2_2}_{\text{Term 1}} \underbrace{- ((x^{(i)})^T u)^2 - ((x^{(i)})^T u)^2}_{\text{Cross Terms}} + \underbrace{((x^{(i)})^T u)^2}_{\text{Last Term}}$$
    
    The last and first term cancel each other out, and we get the following:
    
    $$|x^{(i)}|^2_2 - ((x^{(i)})^T u)^2 - ((x^{(i)})^T u)^2 + ((x^{(i)})^T u)^2$$
    
    
    2) Moving our result from 1) up to multiple examples, we get:
    
    $$J(u) = \sum_{i=1}^{n} \left( |x^{(i)}|^2_2 - ((x^{(i)})^T u)^2 \right)$$
    
    To minimize now this objective, we need to maximize the second subtracted term, so we can say:
    
    $$\arg\min_{u} J(u) = \arg\max_{u} \sum_{i=1}^{n} ((x^{(i)})^T u)^2$$
    
    The covariance matrix of the data is defined as $\Sigma = \sum_{i=1}^{n} x^{(i)}(x^{(i)})^T$. Using this for our problem above, gives us:
    
    $$\sum_{i=1}^{n} ((x^{(i)})^T u)^2 = \sum_{i=1}^{n} (u^T x^{(i)}) ( (x^{(i)})^T u) = u^T \left( \sum_{i=1}^{n} x^{(i)}(x^{(i)})^T \right) u = u^T \Sigma u$$
    
    The unit vector $u$ that maximizes this quadratic form is the eigenvector of the covariance matrix $\Sigma$ with the largest eigenvalue. This eigenvector must be by definition the first principal component of the data.

\end{answer}
